\item En una situación en que los datos se conocen a tres cifras significativas, se escribe $6.379 \text{ m} = 6.38 \text{ m}$ y $6.374 \text{ m} = 6.37 \text{ m}$. Cuando un número termina en $5$, arbitrariamente se elige escribir $6.375 \text{ m} = 6.38 \text{ m}$. Igual se podría escribir $6.375 \text{ m} = 6.37 \text{ m}$, ``redondeando hacia abajo'' en lugar de ``redondear hacia arriba'', porque el número $6.375$ se cambiaría por incrementos iguales en ambos casos. Ahora considere una estimación del orden de magnitud en la cual los factores de cambio, más que los incrementos, son importantes. Se escribe $500 \text{ m} \sim 10^3 \text{ m}$ porque $500$ difiere de $100$ por un factor de $5$, mientras difiere de $1000$ solo por un factor de $2$. Escriba $437 \text{ m} \sim 10^3 \text{ m}$ y $305 \text{ m} \sim 10^2 \text{ m}$. ¿Qué distancia difiere de $100 \text{ m}$ y de $1000 \text{ m}$ por factores iguales de modo que lo mismo se podría escoger representar su orden de magnitud como $\sim 10^2 \text{ m}$ o como $\sim 10^3 \text{ m}$?
